\renewcommand{\part}[1]{\textbf{\large Part \Alph{partCounter}}\stepcounter{partCounter}\\} % part macro
\newcommand{\solution}{\textbf{\large Solution}} % solution macro

% General writing
\newcommand{\bracket}[1]{\left(#1\right)} % for automatic resizing of brackets
\newcommand{\sbracket}[1]{\left[#1\right]} % for automatic resizing of square brackets
\newcommand{\mset}[1]{\left\{#1\right\}} % for automatic resizing of curly brackets
\newcommand{\defeq}{\coloneqq} % the "defined as" command
\newcommand{\ndash}{--\ } % for producing a dash
\newcommand{\mdash}{---\ } % for producing a longer dash
\newcommand{\varemdash}[1][10pt]{%
  \makebox[#1]{\leaders\hbox{---}\hfill\kern0pt}%
}
\newcommand{\halfhrule}[1][0.5]{\begin{center}\varemdash[#1\linewidth]\end{center}}
\newcommand{\by}{\times} % used for talking about m by n matrices
\newcommand{\eg}{e.g.,\ } % for example
\newcommand{\ie}{i.e.,\ } % i.e.
\newcommand{\wrt}{w.r.t.\ } % with respect to 
\newcommand{\RNum}[1]{\uppercase\expandafter{\romannumeral #1\relax}} % uppercase roman numerals
\renewcommand{\cite}{\textcite} % easier citing
\newcommand{\hide}[1]{} % useful for temporarily removing some code

% Macro for variable with label/descriptive (Roman) subscript
\newcommand{\psub}[2]{\ensuremath{#1_{\mathrm{#2}}}}
% Macro for variable with variable (italic) subscript
\newcommand{\vsub}[2]{\ensuremath{#1_{#2}}}

\newcommand{\Reff}{\ensuremath{\psub{R}{e}}}
\newcommand{\ReffGCS}{\ensuremath{\psub{R}{e,GCS}}}
\newcommand{\SN}{\ensuremath{\psub{S}{N}}}
\newcommand{\MV}{\ensuremath{\vsub{M}{V}}}
\newcommand{\Pblue}{\ensuremath{\psub{P}{blue}}}
\newcommand{\Dfb}{\ensuremath{\Delta\psub{f}{b}}}
\newcommand{\meanPblue}{\ensuremath{\langle\psub{P}{blue}\rangle}}
\newcommand{\DeltameanPblue}{\ensuremath{\Delta\langle\psub{P}{blue}\rangle}}
\newcommand{\DeltameanPblueBG}{\ensuremath{\Delta\langle\psub{P}{blue}\rangle_{\rm bgsub}}}
\newcommand{\Rbar}{\ensuremath{\Bar{R}}}
\newcommand{\Rbarblue}{\ensuremath{\psub{\Bar{R}}{blue}}}
\newcommand{\Rbarred}{\ensuremath{\psub{\Bar{R}}{red}}}


\newcommand{\rednote}[1]{{\color{red} #1}} % for red text
\newcommand{\bluenote}[1]{{\color{blue} #1}} % for blue text
\newcommand{\greennote}[1]{{\color{ForestGreen} #1}} % for green text
% \renewcommand{\rednote}[1]{{\color{RoyalPurple} #1}} % alternative to bold text required for AAS review

% setup for underline 
\renewcommand{\ULdepth}{1.8pt}
\contourlength{0.8pt}

\newcommand{\myuline}[1]{%
  \uline{\phantom{#1}}%
  \llap{\contour{white}{#1}}%
}

\newcommand{\asymcomment}[2]{%
	\parbox[t]{#1}{\raggedright #2}%
}

% Blackboard Maths Symbols
\newcommand{\N}{\mathbb{N}} % natural numbers
\newcommand{\Q}{\mathbb{Q}} % rational numbers
% \newcommand{\Z}{\mathbb{Z}} % integers
\newcommand{\R}{\mathbb{R}} % real numbers
\newcommand{\C}{\mathbb{C}} % complex numbers
\newcommand{\E}{\mathbb{E}} % expectaion operators
\renewcommand{\P}{\mathbb{P}} % probability measures

% Calligraphic Maths Symbols
\newcommand{\F}{\mathcal{F}} % sigma-field

% Bold Maths Symbols
\newcommand{\1}{\mathbf{1}} % characteristic function

% Maths operations
\newcommand{\union}{\cup} % union
\newcommand{\intersect}{\cap} % intersection
\newcommand{\directsum}{\oplus} % direct sum
\newcommand{\Union}{\bigcup} % bigger union symbol
\newcommand{\Intersect}{\bigcap} % bigger intersection symbol
\newcommand{\Directsum}{\bigoplus} % bigger direct sum symbol
% \newcommand{\tensor}{\otimes} % tensor
% \newcommand{\Tensor}{\bigotimes} % bigger tensor symbol
\newcommand{\remove}{\setminus} % set difference
\renewcommand{\hom}[2]{\operatorname{Hom}\bracket{{#1}, {#2}}} % Hom sets

% Maths operations
\newcommand{\powerset}[1]{\mathcal{P}\bracket{#1}} % powerset
\renewcommand{\ip}[3]{\left \langle #1, #2 \right \rangle_{#3}} % inner product
\newcommand{\cardinality}[1]{\operatorname{Card}\bracket{#1}} % cardinality
\newcommand{\inv}[1]{#1^{-1}} % inverse
\newcommand{\adj}[1]{#1^{*}} % adjoint
\newcommand{\conj}[1]{\overline{#1}} % conjugate
\newcommand{\vspan}[1]{\operatorname{Span}\left\{#1\right\}} % vector space span
\renewcommand{\rank}[1]{\operatorname{Rank}\bracket{#1}} % rank of a linear transformation
\newcommand{\diag}[1]{\operatorname{Diag}\bracket{#1}} % diagonal elements of a square matrix
\renewcommand{\tr}[1]{\operatorname{Tr}\bracket{#1}} % trace
\newcommand{\kernel}[1]{\operatorname{Ker}\bracket{#1}} % kernel
\newcommand{\im}[1]{\operatorname{Im}\bracket{#1}} % image
\newcommand{\range}[1]{\operatorname{Ran}\bracket{#1}} % range
\newcommand{\col}[1]{\operatorname{Col}\bracket{#1}} % column space
\renewcommand{\exponential}[1]{\operatorname{exp}\bracket{#1}} % exponential function
\newcommand{\floor}[1]{\left\lfloor#1\right\rfloor} % macro for the floor function
\renewcommand{\ceil}[1]{\left\lceil#1\right\rceil} % macro for the ceiling function
\newcommand{\dx}{\,\mathrm{d}x} % differential of x
\newcommand{\dy}{\,\mathrm{d}y} % differential of y
\newcommand{\dt}{\,\mathrm{d}t} % differential of t
\newcommand{\de}[1]{\,\mathrm{d}#1} % general differential 

\newcommand{\partfrac}[2]{{\ifmmode{\ensuremath{\frac{\partial #1}{\partial #2}}}\else$\frac{\partial #1}{\partial #2}$\fi}} % partial derivative fraction
\newcommand{\partfracc}[3]{{\ifmmode{\ensuremath{\left.\partfrac{#1}{#2}\right|_{#3}}}\else$\left.\partfrac{#1}{#2}\right|_{#3}$\fi}} % partial derivative fraction with constants
\newcommand{\partfraccp}[3]{{\ifmmode{\ensuremath{\left(\partfrac{#1}{#2}\right)_{#3}}}\else$\left(\partfrac{#1}{#2}\right)_{#3}$\fi}} % partial derivative fraction with constants in parenthesis
\newcommand{\partfracd}[3]{{\ifmmode{\ensuremath{\left.\partfrac{}{#2}#1\right|_{#3}}}\else$\left.\partfrac{}{#2}#1\right|_{#3}$\fi}} % partial derivative fraction with constants in parenthesis

\newcommand{\partfracbig}[2]{{\ifmmode{\ensuremath{\dfrac{\partial #1}{\partial #2}}}\else$\dfrac{\partial #1}{\partial #2}$\fi}} % partial derivative fraction

% tensor display macros
\newcommand{\tensor}[3]{{\ifmmode{\ensuremath{#1^{#2}\!_{#3}}}\else$#1^{#2}\!_{#3}$\fi}} % tensor parameter, indices up, indices down 
\newcommand{\tensdu}[3]{{\ifmmode{\ensuremath{#1_{#2}\,^{#3}}}\else$#1_{#2}\,^{#3}$\fi}} % tensor parameter, indices down, indices up 
\newcommand{\contra}[2]{{\ifmmode{\ensuremath{#1^{#2}}}\else$#1^{#2}$\fi}} % contravariant tensor parameter, indices up
\newcommand{\covar}[2]{{\ifmmode{\ensuremath{#1_{#2}}}\else$#1_{#2}$\fi}} % covariant (one-form) tensor parameter, indices down

\newcommand{\tenstype}[2]{{\ifmmode{\ensuremath{\begin{pmatrix}#1 \\ #2\end{pmatrix}}}\else$\begin{pmatrix}#1 \\ #2\end{pmatrix}$\fi}} % tensor parameter, indices up, indices down 

% optimisation objectives
\DeclareMathOperator*{\argmin}{argmin} % argument which minimises some criterion
\DeclareMathOperator*{\argmax}{argmax} % argument which maximises some criterion
\DeclareMathOperator*{\arginf}{arginf} % argument that gives the infimum of some criterion
\DeclareMathOperator*{\argsup}{argsup} % argument that gives the supremum of some criterion

% Statistical operations
\newcommand{\given}{\vert} % given (used in conditional probability)
\newcommand{\suchthat}{\;\middle|\;} % such that
\newcommand{\probspace}{\bracket{\Omega, \F, \P}} % a probability space
\newcommand{\borel}[1]{\B\bracket{#1}} % the Borel sigma-field
\newcommand{\prob}[1]{\P\bracket{#1}} % Probability
\newcommand{\cprob}[2]{\P\bracket{#1 \middle\given #2}} % Conditional probability
\newcommand{\expect}[1]{\E\bracket{#1}} % Expectation
\newcommand{\cexpect}[2]{\E\bracket{#1 \middle\given #2}} % Conditional expectation
\renewcommand{\var}[1]{\operatorname{Var}\bracket{#1}} % Variance
\newcommand{\cvar}[2]{\operatorname{Var}\bracket{#1 \middle\given #2}} % Conditional variance
\newcommand{\cov}[2]{\operatorname{Cov}\bracket{#1, #2}} % Covariance
\newcommand{\ccov}[3]{\operatorname{Cov}\bracket{#1, #2 \middle\given #3}} % Conditional covariance
\newcommand{\corr}[2]{\operatorname{Corr}\bracket{#1, #2}} % Correlation
\newcommand{\ccorr}[3]{\operatorname{Corr}\bracket{#1, #2 \middle\given #3}} % Conditional correlation

\newcommand{\nobarfrac}{\genfrac{}{}{0pt}{}}
% Statistical symbols
\newcommand{\indep}{\perp\!\!\!\!\!\perp} % statistical independence
\newcommand{\sigfield}[1]{\sigma\bracket{#1}} % sigma-field
\newcommand{\evalue}[1]{\lambda_{#1}} % eigenvalue

% Include macro for cursive 'r'
% \def\rcurs{{\mbox{$\resizebox{.16in}{.08in}{\includegraphics{ScriptR}}$}}}
% \def\brcurs{{\mbox{$\resizebox{.16in}{.08in}{\includegraphics{BoldR}}$}}}
% \def\hrcurs{{\mbox{$\hat \brcurs$}}}

\def\rcurs{{\mbox{$\resizebox{.11in}{.08in}{\includegraphics[trim= 1em 0 10em 0,clip]{preamble/ScriptR.pdf}}$}}}
\def\brcurs{{\mbox{$\resizebox{.11in}{.08in}{\includegraphics[trim= 1em 0 10em 0,clip]{preamble/BoldR.pdf}}$}}}
\def\hrcurs{{\mbox{$\mathbf{\hat \brcurs}$}}}


\newcommand{\stkout}[1]{\ifmmode\text{\sout{\ensuremath{#1}}}\else\sout{#1}\fi}

\newcommand{\epsO}{{\ifmmode{\ensuremath{\varepsilon_0}}\else$\varepsilon_0$\fi}}

\newcommand{\muO}{{\ifmmode{\ensuremath{\mu_0}}\else$\mu_0$\fi}}

\newcommand{\nab}{{\ifmmode{\ensuremath{\gradient}}\else$\gradient$\fi}}

\newcommand{\del}{{\ifmmode{\ensuremath{\partial}}\else$\partial$\fi}}

\makeatletter
\renewcommand*\env@matrix[1][\arraystretch]{%
  \edef\arraystretch{#1}%
  \hskip -\arraycolsep
  \let\@ifnextchar\new@ifnextchar
  \array{*\c@MaxMatrixCols c}}
\makeatother

% Some greek characters we use as vectors, so need to be bolded.
\newcommand{\vect}[1]{\boldsymbol{\mathbf{#1}}}

\newcommand{\vmu}{{\ifmmode{\ensuremath{\Vec{\mu}}}\else$\Vec{\mu}$\fi}}

% this defines a new command Acolorboxed, to make equation boxes red by default, or color can be specified using e.g.\Acolorboxed[green]{...}
\makeatletter
\newcommand*\Acolorboxed[2][red]{%
   \let\bgroup{\romannumeral-`}%
   \@Acolorboxed{#1}#2&&\ENDDNE
}
\def\@Acolorboxed#1#2&#3&#4\ENDDNE{%
  \ifnum0=`{}\fi
  \setbox\z@\hbox{$\displaystyle#2{}\m@th$\kern\fboxsep \kern\fboxrule}%
  \edef\@tempa{\kern\wd\z@ & \kern-\the\wd\z@ \fboxsep\the\fboxsep \fboxrule\the\fboxrule}%
  \@tempa
  \fcolorbox{#1}{white}{\m@th$\displaystyle#2#3$}%
} 
\makeatother

\newcounter{framecounter}
\newenvironment{myframed}[2][]
    {\refstepcounter{framecounter}% Increment the counter
    \ifx\relax#2\relax % Check if the second argument is empty
    \else
        \label[Box]{#2}% If not, set the label
    \fi
    \mdfsetup{%
    frametitle={%
       {Box \theframecounter: #1}}% Display the counter and title
    }%
    \begin{mdframed}[backgroundcolor=brown!20, linewidth=1.5pt, font=\small]}
    {\end{mdframed}
    }

% \newcommand{\rednote}[1]{{\color{red} #1}} % for red text

\newtcolorbox{mymathbox}[1][]{colback=white, sharp corners, #1}

\newtcolorbox{claimbox}{
  colback=gray!6,
  colframe=gray!45,
  boxrule=0.4pt,
  arc=2pt,
  left=6pt,
  right=6pt,
  top=5pt,
  bottom=5pt,
  before skip=6pt,
  after skip=10pt
}
\newenvironment{claimquote}
{\begin{quote}\itshape\noindent} % \textbf{Claim.}\ 
{\end{quote}}

\newcommand\MBH{\hbox{$M_\text{BH}$}} % Black Hole Mass
\newcommand\Msun{\hbox{ M$_\odot$}} % Solar Mass
\newcommand\Lsun{\hbox{ L$_\odot$}} % Solar luminosity
\newcommand\Rsun{\hbox{ R$_\odot$}} % Solar radius
\newcommand\Mstar{\hbox{ M$_\star$}} % Stellar Mass
\newcommand\arcmin{\mbox{$^\prime$}}% 

% \newcommand{\Comment}[1]{\textcolor{red}{\textit{#1}}}
\DeclareUnicodeCharacter{2609}{\odot}
\DeclareUnicodeCharacter{2605}{\star}
% Astronomy specific units
% common astronomical symbols
\let\sun\odot
\let\bh\bullet
\newcommand*\solarmass{\si{\solarmass}}
\newcommand*\solarluminosity{\si{\solarluminosity}}
% astrophysical units
\DeclareSIUnit\lightspeed{$c$}
\DeclareSIUnit\rydberg{Ry}
\DeclareSIUnit\erg{erg}
\DeclareSIUnit\magnitude{mag}
\DeclareSIUnit\mag{mag}
\DeclareSIUnit\jansky{Jy}
\DeclareSIUnit\gauss{G}
\DeclareSIUnit\h{$h$}
\DeclareSIUnit\hseven{$h$_7}
\DeclareSIUnit\parsec{pc}
\DeclareSIUnit\year{yr}
\DeclareSIUnit\solarluminosity{\ensuremath{L_\sun}}
\DeclareSIUnit\solarmass{\ensuremath{M_\sun}}
\DeclareSIUnit\solarmassinenergy{\ensuremath{M_\sun|c^2}}
\DeclareSIUnit\solarradius{\ensuremath{R_\sun}}
% redeclared units
% \DeclareSIUnit\arcsecond{as}
\DeclareSIUnit\astronomicalunit{AU}
\DeclareSIUnit\clight{\ensuremath c}

\renewcommand{\qedsymbol}{\ensuremath{\blacksquare}}

\definecolor{codegreen}{rgb}{0,0.6,0}
\definecolor{codegray}{rgb}{0.5,0.5,0.5}
\definecolor{codepurple}{rgb}{0.58,0,0.82}
\definecolor{backcolour}{rgb}{0.95,0.95,0.92}

\lstdefinestyle{mystyle}{
    backgroundcolor=\color{backcolour},   
    commentstyle=\color{codegreen},
    keywordstyle=\color{magenta},
    numberstyle=\tiny\color{codegray},
    stringstyle=\color{codepurple},
    basicstyle=\ttfamily\footnotesize,
    breakatwhitespace=false,         
    breaklines=true,                 
    captionpos=b,                    
    keepspaces=true,                 
    numbers=left,                    
    numbersep=5pt,                  
    showspaces=false,                
    showstringspaces=false,
    showtabs=false,                  
    tabsize=2
}

\lstset{style=mystyle}

\setlength\bibindent{2em}

% \makeatletter
% \renewcommand\NAT@bibsetnum[1]{\settowidth\labelwidth{\@biblabel{#1}}%
%    \setlength{\leftmargin}{\bibindent}\addtolength{\leftmargin}{\dimexpr\labelwidth+\labelsep\relax}%
%    \setlength{\itemindent}{-\bibindent}%
%    \setlength{\listparindent}{\itemindent}
% \setlength{\itemsep}{\bibsep}\setlength{\parsep}{\z@}%
%    \ifNAT@openbib
%      \addtolength{\leftmargin}{\bibindent}%
%      \setlength{\itemindent}{-\bibindent}%
%      \setlength{\listparindent}{\itemindent}%
%      \setlength{\parsep}{0pt}%
%    \fi
% }
% \makeatother

%\newtotcounter{citesnum}
\def\oldcite{} \let\oldcite=\cite
\def\cite{\stepcounter{citesnum}\oldcite}

%\newtotcounter{bibnum}
\def\oldbibitem{} \let\oldbibitem=\bibitem
\def\bibitem{\stepcounter{bibnum}\oldbibitem}

%Now using biblatex, so the bibnum counter is no longer correct. Following code changes this to add references printed, cited and total cited
\usepackage{totcount}
\newtotcounter{bibprintedtotal}   % total printed across all bibliographies
\newcounter{bibprinted}           % printed in the *current* \printbibliography

\AtBeginBibliography{%
  \setcounter{bibprinted}{0}%
}

\AtEveryBibitem{%
  \stepcounter{bibprinted}%
  \stepcounter{bibnum}%
}

% Count citation occurrences (including repeats)
\newtotcounter{citehits}      % counts each cited key occurrence (incl repeats)
\newtotcounter{citedunique}   % counts unique keys cited at least once
\newtotcounter{citerepeats}   % counts repeat citations of already-seen keys

% \AtEveryCitekey{%
%   % count a "hit" for every citekey, every time
%   \stepcounter{citehits}%

%   % stable key
%   \edef\blx@tmpkey{\thefield{entrykey}}%

%   % count uniques vs repeats
%   \ifcsundef{blx@seen@\blx@tmpkey}
%     {%
%       \csdef{blx@seen@\blx@tmpkey}{1}%
%       \stepcounter{citedunique}%
%     }
%     {%
%       \stepcounter{citerepeats}%
%     }%
% }


\newcommand{\chapteropen}{%
  \clearpage
  % \ifodd\value{page}\else   % uncomment these lines for chapters starting on odd pages
  %   \null
  %   \thispagestyle{empty}%
  %   \newpage
  % \fi
}
    
\renewcommand{\bibname}{BIBLIOGRAPHY}

\UseTblrLibrary{functional}
%% Commands for Typesetting stuff in table
\NewExpandableDocumentCommand{\WP}{m m}{\textbf{WP-#1: #2}}
\NewExpandableDocumentCommand{\Dx}{O{1} m}
{\textbf{D\textsubscript{#1}:} \textit{#2}}
\NewExpandableDocumentCommand{\Mx}{O{1} m}
{\textbf{M\textsubscript{#1}:} \textit{#2}}
%% Functional for coloring cells
\IgnoreSpacesOn
\definecolor{mycolor}{HTML}{91c1eb}
\newcommand{\mystr}{x}
\prgNewFunction \qColor {} {
    \intStepOneInline {3} {\arabic{rowcount}} {
        \intStepOneInline {3} {\arabic{colcount}} {
            \strSet \lTmpaStr {\cellGetText {##1} {####1} }
            \strVarIfEmptyF \lTmpaStr {
                \cellSetStyle{##1}{####1}{bg=mycolor}
                \cellSetText{##1}{####1}{}
            }
        }   
    }
}
\IgnoreSpacesOff

\newcommand{\settitles}{
    \renewcommand{\contentsname}{TABLE OF CONTENTS\vspace{-0.20in}}
    \renewcommand\listfigurename{LIST OF FIGURES\vspace{-0.20in}}
    \renewcommand\listtablename{LIST OF TABLES\vspace{-0.20in}}
    \renewcommand{\bibname}{REFERENCES}
    \renewcommand{\chaptername}{CHAPTER}
}

\setlength{\abovecaptionskip}{6pt}
\setlength{\belowcaptionskip}{6pt}  % adjust for caption-to-table spacing
\renewcommand{\arraystretch}{1.2}  % default is 1.0 for spacing in tables

\setlength{\heavyrulewidth}{0.04em}
\setlength{\lightrulewidth}{0.02em}
\setlength{\cmidrulewidth}{0.03em}
\newcommand{\doubletoprule}{%
\toprule
\addlinespace[0.15em]
\toprule
}

\newcommand{\tablecomments}[1]{%
  \vspace{0.5em}  
  \par\noindent
  \begin{minipage}{\linewidth}
  \setstretch{1.0}
  \footnotesize
  \textsc{Note}---#1
  \end{minipage}\par
}
% Define the environment with a helper command for the attribution
\newcommand{\bqauthor}{}
\newenvironment{bquote}[1]
  {\renewcommand{\bqauthor}{#1}%
   \begin{adjustwidth}{1in}{\dimexpr\paperwidth-7.5in}%
   \itshape\justifying}
  {%
   \begin{flushright}
     \textbf{\itshape--- \bqauthor}
   \end{flushright}
   \end{adjustwidth}
}


% --- Widow & orphan control ---
\clubpenalty=10000
\widowpenalty=10000
\displaywidowpenalty=10000

% Give LaTeX flexibility to avoid bad breaks
\emergencystretch=2em

\newcommand{\references}{
    \if@openright\cleardoublepage\else\clearpage\fi
    \phantomsection
    \addcontentsline{toc}{chapter}{\bibname} % 
    \printbibliography[title={\bibname}]
}


% Reproduce AASTeX-style \dataset and \software commands
\makeatletter

% Save existing \doi if it exists
\let\savedoi\doi

% \dataset{URL}
% \dataset[Label]{URL}
% Also allows \dataset{\doi{10.xxxx/yyyy}}
\newcommand{\dataset}{%
  \begingroup
  \def\doi##1{https://doi.org/##1}%
  \@ifnextchar[{\ydataset}{\xdataset}%
}

\newcommand{\xdataset}[1]{%
  \href{#1}{[DATASET]}%
  \endgroup
}

\newcommand{\ydataset}[2][]{%
  \href{#2}{#1}%
  \endgroup
}

% \software{Astropy, NumPy, ...}
\newcommand{\software}[1]{%
  \par\vskip 6pt%
  {\large\itshape Software: }#1\par
}

\makeatother
